problem:  
Let \((M^n, g)\) be a complete, noncompact Riemannian manifold with nonnegative Ricci curvature and **maximal (Euclidean) volume growth**, i.e.  
\[
\lim_{r\to\infty} \frac{\operatorname{Vol}(B(p,r))}{\omega_n r^n} = \theta(M,g) > 0.
\]
By results of Cheeger–Colding, any blow-down sequence \((M, r_i^{-2}g, p)\) with \(r_i \to \infty\) admits a subsequence converging in the pointed Gromov–Hausdorff sense to a metric cone \(C(X)\) over a compact length space \(X\), where \(X\) satisfies \(\operatorname{RCD}^*(n-2,n-1)\) conditions and \((C(X),d_C)\) is noncollapsed.

**Problem:**  
Assume, in addition, that such a tangent cone at infinity is *unique* (i.e. independent of the sequence \(r_i\to\infty\)). Under these hypotheses:
> Must the cross-section \(X\) necessarily be a smooth Einstein manifold satisfying \(\operatorname{Ric}_X = (n-2)h\) (with \(h\) its metric), equivalently implying that the asymptotic cone \(C(X)\) is a smooth Ricci-flat Riemannian manifold outside the vertex?

In other words, does uniqueness of the asymptotic cone plus maximal volume growth enforce *analytic and geometric smoothness* of the metric cone’s cross-section in the nonnegative Ricci setting?

Why is it a "good" problem:  
This problem refines the classical Cheeger–Colding framework by coupling volume growth rigidity with uniqueness of tangent cones, a condition intimately connected to the global geometry at infinity. While theory guarantees the existence of metric cone limits and describes partial regularity (smooth regular part, singular set of codimension ≥ 2), it remains unknown whether **uniqueness** of the asymptotic cone rules out such singularities altogether.  

The question elegantly probes the boundary between **metric-measure geometry** and **smooth Riemannian geometry**. A positive answer would establish a new rigidity theorem connecting asymptotic volume growth, tangent cone uniqueness, and Einstein smoothness—while a counterexample would reveal deep structural subtleties in the analytic limits of Ricci ≥ 0 manifolds.  

Its formulation is minimal (no extraneous assumptions, clearly geometric, directly founded on existing theory), yet its resolution would require genuinely new insight into analytic regularity of noncollapsed Ricci limit spaces, potentially reshaping the understanding of geometric stability and smoothness in Ricci limit theory.